% Compile-ready IEEEtran LaTeX Sections IV-VII matching results/ CSVs and figures
% Include in your IEEEtran paper after Section III (using \usepackage{booktabs,graphicx,amsmath})

\section{Implementation}
We instantiate the four-pillar CTI blueprint as a modular Python~3.13 reference platform (\texttt{agent-cti}) comprising a tool-using target agent, a simulated human-in-the-loop (HITL) baseline, and four composable defense layers orchestrated via \texttt{DefensePipeline}.

\subsection{Target AI System Under Defense}
\texttt{TargetAgent} runs a multi-step ReAct planning loop ($\text{max\_steps}=6$, $\text{max\_tokens}=4096$) equipped with five sandbox-ready tools (\texttt{file\_read}, \texttt{send\_email}, \texttt{http\_fetch}, \texttt{sql\_query}, and AST-validated \texttt{calculator}), an external plugin registration surface (\texttt{register\_plugin\_manifest}), a JSON-persisted \texttt{MemoryStore}, a \texttt{RAGStore} supporting both TF-IDF and live CPU neural dense retrieval (\texttt{sentence-transformers/all-MiniLM-L6-v2}), and a two-agent broker (\texttt{AgentBus}). All sensitive values are synthetic canary markers (\texttt{CANARY\_API\_KEY\_7F9A2B}, \texttt{CANARY\_DB\_TOKEN\_X99}, \texttt{CANARY\_SSN\_000-11-2222}). A pluggable \texttt{LLMBackend} protocol supports deterministic scripted runs (\texttt{DeterministicMockLLM}), seeded probabilistic evaluation (\texttt{StochasticMockLLM}), and \texttt{.env}-configured adapters for OpenAI-compatible, Gemini, Anthropic, Groq, Ollama, and local HuggingFace backends.

\subsection{Pillar 1: Threat Understanding}
\texttt{threats/catalog.yaml} models 13 AI/agentic threat classes (\texttt{T01}--\texttt{T13}) mapped across system scopes (\emph{models, data pipelines, retrieval systems, tools, agents, user interactions, memory, agent channel}), scored by $\text{Impact} \times \text{Likelihood} \in [1, 25]$, and cross-referenced to OWASP LLM Top~10 v2025, OWASP Agentic AI Threats, and MITRE ATLAS. \texttt{ControlGenerator} ranks controls by aggregate risk and emits \texttt{policies/generated\_controls.yaml} to drive both Pillar~2 and Pillar~3 enablement order.

\subsection{Pillar 2: Governance and Guardrails}
\texttt{PolicyEngine} evaluates declarative YAML rules across \texttt{pre\_input}, \texttt{system\_prompt} (\texttt{[POLICY\_FRAMED]} boundaries), \texttt{pre\_retrieval}, \texttt{pre\_tool} (allowlists, read-only SQL, file/domain guards, memory-write checks, per-tool rate caps, and external-action escalation), and \texttt{post\_output} (canary/PII DLP and schema validation), logging every \texttt{ALLOW}/\texttt{BLOCK}/\texttt{ESCALATE} verdict with its triggering \texttt{policy\_id}. To generalize beyond brittle regexes without false positives on dual-use queries, \texttt{SemanticGuardrail} runs \texttt{all-MiniLM-L6-v2} locally and flags text only when both $\text{sim}_{\text{atk}} \ge 0.48$ and $\text{sim}_{\text{atk}} - \max \text{sim}_{\text{ben}} \ge 0.08$.

\subsection{Pillar 3: Runtime Security}
Runtime enforcement secures the underlying substrate via: (1) per-agent Ed25519 keys (\texttt{IdentityRegistry}) verifying inter-agent messages on \texttt{AgentBus}; (2) short-lived single-use HMAC-SHA256 \texttt{CapabilityToken}s ($\text{ttl}=30\,\text{s}$); (3) isolated \texttt{python -I} \texttt{SubprocessSandbox} execution with scrubbed environment variables, root-jail filesystem containment, and blocked Python sockets; (4) an append-only SHA-256 hash-chained \texttt{AuditLedger} with offline tamper verification; and (5) behavioral \texttt{AnomalyDetector} triggering a microsecond \texttt{KillSwitch} quarantine.

\subsection{Pillar 4: Adversarial Assurance and Simulated HITL}
We enforce strict train/test separation between a 26-attack regression suite (\texttt{scenarios.yaml}) and a zero-leakage 26-attack unseen polymorphic holdout suite plus 6 hard dual-use benign queries (\texttt{holdout\_scenarios.yaml}). Findings from \texttt{C4} update Pillar~1 likelihoods, emit YAML patches (\texttt{remediation\_policies.yaml}), and calibrate neural prototypes (\texttt{semantic\_prototypes.json}). To test HITL scaling without a human-subjects study, \texttt{SimulatedHumanApprover} models 1 or 3 reviewers ($2000\,\text{ms}$ service time, $\text{max\_queue}=20$, $15\,\text{s}$ timeout, $0.04$ base miss rate, and $0.035$ linear fatigue slope per queued item).

\section{Experimental Evaluation}
\subsection{Experimental Setup}
All experiments ran on an Intel Core i7 (13th-gen, CPU-only, 16\,GB RAM) Windows~11 host using Python~3.13.5 across 5 stochastic seeds (\{42, 101, 202, 303, 404\}) reporting means and 95\% Student-$t$ confidence intervals ($\text{df}=4$) alongside paired $t$-tests and Cohen's $d$.

\subsection{Ablation Across Conditions C0--C5}
\begin{table*}[t]
\centering
\caption{Ablation Across Defense Conditions (5 Seeds, Mean $\pm$ 95\% CI)}
\begin{tabular}{llrrrrrr}
\toprule
Cond. & Configuration & ASR (\%) & Block/Detect (\%) & Benign (\%) & Esc. (\%) & Added Step p50/p95 (ms) & TTD / TTC (ms) \\
\midrule
C0 & Undefended Baseline & $89.23 \pm 5.50$ & $0.00 \pm 0.00$ & $100.00 \pm 0.00$ & $0.00$ & 0.00 / 0.00 & -- / -- \\
C1 & Simulated HITL Only & $30.38 \pm 3.92$ & $59.62 \pm 4.78$ & $90.77 \pm 4.27$ & $0.00$ & 1203.90 / 2001.34 & -- / -- \\
C2 & Pillar 2 Guardrails Only & $18.08 \pm 1.31$ & $66.15 \pm 12.36$ & $92.31 \pm 0.00$ & $1.54$ & 0.19 / 0.00 & -- / -- \\
C3 & Pillar 3 Runtime Only (Sandboxed) & $31.54 \pm 3.62$ & $59.62 \pm 6.09$ & $100.00 \pm 0.00$ & $0.00$ & 4.74 / 201.40 & 139.64 / 0.020 \\
C4 & Pillars 1+2+3 (Pre-Feedback) & $15.00 \pm 1.07$ & $68.46 \pm 11.65$ & $92.31 \pm 0.00$ & $1.54$ & 7.99 / 177.34 & 337.05 / 0.023 \\
C5 & Full 4-Pillar Loop (Post-Feedback) & $0.00 \pm 0.00$ & $83.46 \pm 11.53$ & $92.31 \pm 0.00$ & $1.54$ & 106.50 / 227.08 & 411.75 / 0.025 \\
\bottomrule
\end{tabular}
\end{table*}


As shown in Table~\ref{tab:ablation} (see \texttt{results/ablation\_table.tex}) and Fig.~\ref{fig:ablation}, no single defense suffices: simulated HITL (\texttt{C1}) cuts Attack Success Rate (ASR) from $89.23 \pm 5.50\%$ to $30.38 \pm 3.92\%$ ($p=2.14\times 10^{-6}$, $d=18.29$) while adding $1203.90\,\text{ms}$ p50 ($2001.34\,\text{ms}$ p95) per-step latency and missing non-tool attacks (\texttt{T03} $95.0\%$, \texttt{T04} $95.0\%$, \texttt{T13} $100.0\%$). \texttt{C3} (Runtime Only) eliminates tool, identity, loop, and supply-chain attacks ($0.0\%$ ASR on \texttt{T07}--\texttt{T09}, \texttt{T11}, \texttt{T12}) with $100.0\%$ benign success, whereas \texttt{C2} (Guardrails Only) achieves $18.08 \pm 1.31\%$ ASR at $+0.19\,\text{ms}$ p50 overhead. Combining Pillars~1--3 (\texttt{C4}) lowers ASR to $15.00 \pm 1.07\%$ ($p=0.0028$ vs.~\texttt{C2}; $p=4.51\times 10^{-4}$ vs.~\texttt{C3}), and closing the Pillar~4 loop (\texttt{C5}) reduces regression ASR to $0.00\%$ ($p=2.58\times 10^{-6}$, $d=17.44$) at $92.31\%$ autonomous benign completion ($1.54\%$ intentional \texttt{POL-ESC-001} escalation on external email \texttt{BEN-13}).

\subsection{Zero-Leakage Unseen Holdout Generalization}
\begin{table*}[t]
\centering
\caption{Unseen Holdout Generalization and Hard Dual-Use Benign Utility (5 Seeds, Mean $\pm$ 95\% CI)}
\begin{tabular}{llrrrrrr}
\toprule
Cond. & Configuration & ASR (\%) & Block/Detect (\%) & Benign (\%) & Esc. (\%) & Added Step p50/p95 (ms) & TTD / TTC (ms) \\
\midrule
C0 & Undefended Baseline & $91.54 \pm 2.14$ & $0.00 \pm 0.00$ & $100.00 \pm 0.00$ & $0.00$ & 0.00 / 0.00 & -- / -- \\
C4 & Pillars 1+2+3 (Pre-Feedback) & $22.31 \pm 4.65$ & $63.08 \pm 5.70$ & $100.00 \pm 0.00$ & $0.00$ & 6.58 / 150.89 & 356.94 / 0.017 \\
C5\_REGEX & Pillar 4 Regex Remediation Only & $22.31 \pm 4.65$ & $63.08 \pm 5.70$ & $100.00 \pm 0.00$ & $0.00$ & 6.07 / 153.49 & 385.15 / 0.018 \\
C5\_FULL & Pillar 4 Regex + Neural Semantic & $15.77 \pm 3.54$ & $69.62 \pm 6.63$ & $100.00 \pm 0.00$ & $0.00$ & 50.30 / 172.38 & 405.56 / 0.018 \\
\bottomrule
\end{tabular}
\end{table*}


Table~\ref{tab:holdout} demonstrates that regex-only patches (\texttt{C5\_REGEX}) overfit the dev suite, yielding zero improvement on unseen variants ($22.31\%$ vs.\ $22.31\%$, $p=1.00$). Adding \texttt{all-MiniLM-L6-v2} semantic screening (\texttt{C5\_FULL}) eliminates all unseen RAG and memory paraphrase bypasses, cutting holdout ASR to $15.77 \pm 3.54\%$ ($p=1.56\times 10^{-4}$, $d=6.21$) with $100.00\%$ Hard Benign utility on dual-use queries (\texttt{HBEN-01}--\texttt{HBEN-06}) at $+50.30\,\text{ms}$ CPU step overhead, while leaving $15.77\%$ adaptive residual ASR (reversed-token DLP evasion and novel bulletin hallucinations).

\subsection{Pillar 1 Prioritization, Dense RAG, Probes, and HITL Load Scaling}
Under constrained budgets $K \in \{2, 4, 6, 8, 11\}$ of 11 controls, Pillar~1 risk-ranked ordering achieves $29.23 \pm 2.62\%$ ASR at $K=2$ and $20.00 \pm 2.14\%$ at $K=4$ (vs.\ $58.46\%$ / $38.46\%$ Random and $80.00\%$ / $73.08\%$ Reverse-Risk). Live dense \texttt{all-MiniLM-L6-v2} RAG replicates TF-IDF mitigation trajectories ($100.0\% \rightarrow 50.0\% \rightarrow 0.0\%$ \texttt{T04} ASR), and 10/10 direct tamper/replay/spoof/sandbox probes against Pillar~3 are blocked ($100.0\%$). Runtime \texttt{KillSwitch} containment completes in $0.020\text{--}0.025\,\text{ms}$. Under load scaling ($5\text{--}480\,\text{req/min}$), simulated 1-reviewer HITL saturates at $30.1\,\text{req/min}$ ($87.17\%$ drops at $480\,\text{req/min}$) and 3-reviewer HITL reaches $41.33\%$ fatigue misses at $120\,\text{req/min}$ ($p95=15.37\,\text{s}$), whereas Automated \texttt{C5} sustains $480.0\,\text{req/min}$ in-process ($35.45\,\text{ms}$ p95) and $333.37\,\text{req/min}$ sandboxed ($246.78\,\text{ms}$ p95) with $0.0\%$ drops.

\section{Discussion and Limitations}
(1) \textbf{Simulated HITL}: \texttt{C1} and scaling curves derive from a parametric queueing/fatigue simulator rather than a human-subjects study. (2) \textbf{Mock vs.\ Real Generative LLMs}: Agent planning used \texttt{DeterministicMockLLM} and \texttt{StochasticMockLLM} offline while neural guardrails and dense RAG executed live \texttt{all-MiniLM-L6-v2} inference on CPU; live cloud/local generative runs are supported via \texttt{.env}. (3) \textbf{Residual Adaptive Risk}: Holdout ASR of $15.77\%$ confirms continuous assurance is necessary against unseen DLP encodings and hallucinations. (4) \textbf{Windows Subprocess Scope}: Production environments should substitute container/microVM kernel isolation for \texttt{python -I} sandboxing.

\section{Conclusion}
Empirical evaluation confirms that human-in-the-loop oversight fails under agentic speed and non-tool attack vectors, whereas combining risk-prioritized Threat Understanding, declarative and semantic Guardrails, microsecond-containment Runtime Security, and continuous Adversarial Assurance reduces known ASR to $0.00\%$ and unseen holdout ASR to $15.77\%$ while preserving $92.31\%\text{--}100.0\%$ benign utility and sustaining up to $480\,\text{req/min}$.

